% ===== tr-06 块2：第 6 章 pp.22--25 + p.26 参考文献前正文（原文页 22--26）=====
% [边界] p.26 页首是命题 6.6.16 的 (1)(2) 条及其结论（正文延伸入参考文献起始页），其后为 References；参考文献由 Lead 脚本转录，本文件不写。
% [原文 p.22]
于是由定理 6.6.1 可得配对同构 $(\check Y_1,\check D_1)\cong(\check Y_2,\check D_2)$。\footnote{注意：由于 $\check X_i$ 的 Kähler 锥比从 $\check Y_i$ 限制而来的锥高一维，要论证该同构确实保持 Kähler 类还需要一些处理。}

一旦配对 $(\check Y,\check D)$ 由 Looijenga 对的 Torelli 定理确定，定理 6.5.2 便保证 Kähler 类 $[\check\omega]$ 将在沿全纯截面作平移（相对于某个参考截面）的意义下唯一决定 $\check\omega$。注意沿全纯截面作平移不改变 $\check\Omega$。于是我们得到 ALH$^*$ 引力瞬子的 Torelli 定理。

\begin{theorem}{6.6.11}[{[12, 定理 3.9]}]
设 $(X_i,\omega_i,\Omega_i)$ 是 ALH$^*$ 引力瞬子，且存在微分同胚 $F:X_2\cong X_1$ 使得
\[
F^*[\omega_1]=[\omega_2]\in H^2(X_2,\R),\qquad F^*[\Omega_1]=[\Omega_2]\in H^2(X_2,\C).
\]
则存在微分同胚 $f:X_2\to X_1$ 使得 $f^*\omega_1=\omega_2$ 且 $f^*\Omega_1=\Omega_2$。
\end{theorem}

这个 Torelli 定理有几个值得提及的推论。首先，我们在二维情形下部分回答了 Tian--Yau [62] 提出的一个问题。

\begin{corollary}{6.6.12}
设 $Y$ 是弱 del Pezzo 曲面，$D$ 是光滑反典范除子，$X=Y\setminus D$。固定 $Y$ 上一个沿 $D$ 有单极点的亚纯 2-形式 $\Omega$。对 $X$ 的每个 Kähler 类，存在该 Kähler 类中唯一的、具有 $L^2$ 曲率的超 Kähler 度量 $\omega$。
\end{corollary}

\begin{proof}
存在性由 [62, 定理 5.1] 保证。反之，若存在这样一个超 Kähler 度量 $\omega$，则由 [60]，$(X,\omega,\Omega)$ 是引力瞬子。由于引力瞬子的可能微分同胚型只有有限多个，$(X,\omega,\Omega)$ 必为 ALH$^*$ 引力瞬子。于是该推论由定理 6.6.11 推出。\proofend
\end{proof}

在定理 6.2.3 及其记号下，给定 $H_1(D,\Z)$ 中任一本原类，$X$ 中就有相应的 SYZ 纤维化。配对 $(Y,D)$ 的模空间的单值作用在 $H_2(X,\Z)$ 上，特别地诱导出在 $\operatorname{Im}(H_1(D,\Z)\to H_2(X,\Z))\cong\Z^2$ 上的作用。Torelli 定理的一个直接应用是回答 [34] 中提出的一个问题。

\begin{proposition}{6.6.13}[{Lau--Lee-L., [49]}]
对 $\operatorname{Aut}(H_2(X,\Z))$ 中任一单值变换，存在 $X$ 的一个等距，使其在 $H_2(X,\Z)$ 上诱导的映射实现给定的单值变换，并把特殊拉格朗日纤维化送到特殊拉格朗日纤维化。特别地，不同 SYZ 纤维化的辛仿射结构与复仿射结构彼此同构。
\end{proposition}

最后一个应用是 ALH$^*$ 引力瞬子的另一个单值化定理（属于 Hein--Sun--Viaclovsky--Zhang）的另一种证明。
% [原文 p.23]
\begin{theorem}{6.6.14}[{Hein--Sun--Viaclovsky--Zhang, [37]}]
任一 ALH$^*$ 引力瞬子 $(X,\omega,\Omega)$ 都可以通过添加一个光滑反典范除子 $D$ 而紧化为弱 del Pezzo 曲面 $Y$，使得 $\Omega$ 是 $Y$ 上沿 $D$ 有单极点的亚纯 2-形式。
\end{theorem}

在给出证明之前，先对不同几何的超 Kähler 三元组模空间作一次维数核对，是个不错的合理性检验。一方面，带 $I_d$-纤维的有理椭圆曲面的模空间复维数为 $9-d$，沿给定纤维有单极点的亚纯 2-形式在 $\C^*$ 伸缩下唯一，补集的 Kähler 锥实维数为 $11-d$，故总实维数为 $31-3d$。另一方面，带光滑反典范除子的弱 del Pezzo 曲面配对的模空间复维数为 $10-d$，沿给定光滑反典范除子有单极点的亚纯体积形式在 $\C^*$ 伸缩下唯一，补集的 Kähler 锥实维数为 $9-d$，\footnote{这是由于类似于 (6.3.1) 的配对的长正合序列，以及 Calabi ansatz 在无穷远处是恰当的。} 故总实维数又是 $31-3d$。

\begin{proof}[定理 6.6.14 的证明]
考虑由 5-元组 $(Y,D,[\omega],\alpha,c)$ 构成的空间 $\mathcal{A}$，其中
\begin{enumerate}
\item $Y$ 是 $Y_0$ 的形变族中次数为 $d$ 的弱 del Pezzo 曲面；
\item $D\in|-K_Y|$ 是光滑反典范除子；
\item $[\omega]\in H^2(Y\setminus D,\R)$；
\item $\alpha\in H_1(D,\Z)$ 是一个本原类；
\item $c\in\R_+$。
\end{enumerate}
若存在配对同构 $f:(Y,D)\to(Y',D')$ 使得 $f_*\alpha=\alpha'$、$f^*[\omega']=[\omega]$、$c=c'$，则等同两个元素 $(Y,D,[\omega],\alpha,c)$ 与 $(Y',D',[\omega'],\alpha',c)$。$\mathcal{A}$ 上有自然拓扑。类似地，考虑由形如 $(\check Y,\check D,[\check\omega],\check c)$ 的 4-元组构成的空间 $\mathcal{B}$，其中
\begin{enumerate}
\item $\check Y$ 是有理椭圆曲面；
\item $\check D$ 是 $\check Y$ 的 $I_d$-纤维；
\item $[\check\omega]\in H^2(\check Y\setminus\check D,\R)$ 是 Kähler 类；
\item $\check c\in\R_+$。
\end{enumerate}
我们类似地等同 $\mathcal{B}$ 中的元素，$\mathcal{B}$ 上也有自然拓扑。

给定 4-元组 $(\check Y,\check D,[\check\omega],\check c)$，存在 $\check Y$ 上沿 $\check D$ 有单极点的唯一亚纯 2-形式 $\check\Omega$，归一化使其在坏轮上的积分为 $\check c$。对足够大的 $\check c$，由 Hein 的构造存在上同调类 $[\check\omega]$ 中渐近于 $\check c\check\omega_{sf}$ 的 Ricci 平坦度量 $\check\omega$，并且由定理 6.5.2 在平移意义下唯一。特别地，$2\check\omega^2=\check\Omega\wedge\bar{\check\Omega}$。由下面的证明将会看到，“$\check c$ 足够大”这一限制可以去掉。
% [原文 p.24]
给定元素 $(Y,D,[\omega],\alpha,c)\in\mathcal{A}$，由 [62, 定理 5.1] 存在 Tian--Yau 度量 $\omega$ 的唯一伸缩，使其渐近于 $c\omega_C$。由定理 6.2.3，$X$ 中存在纤维类对应于 $\alpha$ 的 SYZ 纤维化。$X$ 经 (6.2.2) 的超 Kähler 旋转可以通过添加 $I_d$-纤维 $\check D$ 而紧化为有理椭圆曲面 $\check Y$，其 Ricci 平坦度量 $\check\omega$ 渐近于某个（可能非标准的）半平坦度量的某个伸缩。总之，存在连续映射
\[
\mathcal{HK}:\mathcal{A}\to\mathcal{B}.
\]
我们指出，经过超 Kähler 旋转后，$H_2(X,\Z)$ 中对应于 $\beta\in H_1(D,\Z)$、满足 $\langle\alpha,\beta\rangle=1$ 的同调类，恰好成为 $(\check Y,\check D)$ 的坏轮（差一个纤维的倍数）。这一信息有助于归一化 $(2,0)$-形式以应用定理 6.6.1。

我们断言，超 Kähler 旋转映射 $\mathcal{HK}$ 是满射。我们先假定该断言来证明定理。取一个 ALH$^*$ 引力瞬子 $(X,\omega,\Omega)$ 以及特殊拉格朗日同调类 $[L]\in H_2(X,\Z)$ 的一个选取。通过调整 $\Omega$ 的相位，可设 $\int_{[L]}\Omega>0$。于是定理 6.6.10 给出 $\mathcal{B}$ 中一个元素，它与 $\mathcal{HK}(Y',D',\alpha',\beta',c')$ 相合。特别地，有双全纯
\[
\check X\cong(X,J),\qquad \check X\cong(X',J'),
\]
其中 $J$（以及 $J'$）表示由 2-形式 $\omega-i\operatorname{Im}\Omega$（以及相应地 $\omega'-\operatorname{Im}\Omega'$）所决定的复结构。注意 $\check Y$ 上沿 $\check D$ 有单极点的亚纯 2-形式在 $\C^*$ 伸缩下唯一。利用上一段的归一化与固定的伸缩 $\check c$，由满射假设下的双全纯有 $[\Omega]=[\Omega']$ 与 $[\omega]=[\omega']$。于是由定理 6.6.11，我们得到引力瞬子的等距 $(X',\omega',\Omega')\cong(X,\omega,\Omega)$。特别地，$X$ 可以紧化为一个与 $Y'$ 双全纯的紧复曲面。

为证明该断言，只需证明 $\mathcal{HK}$ 的像既开又闭。我们先证开性。设 $\mathcal{HK}(Y,D,[\omega],\alpha,c)=(\check Y,\check D,[\check\omega],\check c)$。回忆 $j$-不变量在带光滑反典范除子的弱 del Pezzo 曲面配对的模空间上定义了一个非常值全纯函数。特别地，它是开映射。结合 [19, 命题 3.12]，我们可以通过形变配对 $(Y,D)$（即形变 $D$ 以及形变 $\PP^2$ 上的爆破中心）来实现 $\Omega$ 的所有可能的邻近周期。由 [62, 定理 5.1]，在 $X'$ 上任一充分接近 $[\omega]$ 的 Kähler 类 $[\omega']$ 中，存在渐近于 Calabi ansatz 的 Ricci 平坦度量，\footnote{这里我们通过形变族中的平行移动来认同上同调类。} 其中 $X'=Y'\setminus D'$ 且 $(Y',D')$ 是 $(Y,D)$ 的形变。这意味着经过超 Kähler 旋转后，可以实现 $[\check\Omega]$ 与 $[\check\omega]$ 的周期的所有小形变，其中 $\check\Omega$ 是 $\check Y$ 上沿 $\check D$ 有单极点的亚纯 2-形式。因此由定理 6.6.1，$\mathcal{HK}$ 是开的。最后，我们证明 $\mathcal{HK}$ 是闭的。设在
% [原文 p.25]
$\mathcal{B}$ 中存在收敛序列
\begin{equation}
\mathcal{HK}(Y_i,D_i,[\omega_i],\alpha_i,c_i)=(\check Y_i,\check D_i,[\check\omega_i],\check c_i)\to(\check Y,\check D,[\check\omega],\check c). \tag{6.6.4}
\end{equation}
注意 $D_i$ 的模可以由 $[\check\omega_i-i\operatorname{Im}\check\Omega_i]$ 在椭圆纤维以及 $\check Y_i$ 的坏轮上的积分恢复。由收敛性 (6.6.4)，我们有 $D_i\to D$，其中 $D$ 是某个光滑椭圆曲线。周期 $[\check\omega_i-i\operatorname{Im}\check\Omega_i]$ 在 $H_2(\check X_i,\Z)$ 上的收敛蕴涵着存在弱 del Pezzo 曲面 $Y$ 使得 $(Y_i,D_i)\to(Y,D)$（由 [53, 定理 6.4]）。于是 $[\omega]=[\operatorname{Re}\check\Omega]$（差平行移动），$\alpha$ 由 $\alpha_i$ 经平行移动决定，且 $c=\lim_i c_i$。不难看出 $\mathcal{HK}(Y,D,[\omega],\alpha,c)=(\check Y,\check D,[\check\omega],\check c)$，故 $\mathcal{HK}$ 是闭的。\proofend
\end{proof}
% [Tutor Note: 上一段“有双全纯 $\check X\cong(X,J)$、$\check X\cong(X',J')$”原文两处都写 $\check X$（页 24 可核对），按论证应为 $\check X\cong(X,J)$ 与 $\check X'\cong(X',J')$；又括号内第二个复结构印作 $\omega'-\operatorname{Im}\Omega'$，比第一个少了虚数单位 $i$（页 24 页图可辨）。两处均照录原文，此注给出应有形式。]

由 Torelli 定理自然产生几个问题。回忆给定 ALH$^*$ 引力瞬子 $X$ 以及到某个固定参考 ALH$^*$ 引力瞬子 $X_0$ 的微分同胚，$X$ 的超 Kähler 三元组经该微分同胚诱导出 $H^2(X_0,\R)^3$ 中一个元素。从 ALH$^*$ 引力瞬子空间到 $H^2(X_0,\R)^3$ 的映射是 K3 曲面周期映射的类比，因此 $H^2(X_0,\R^3)$ 是周期区域的类比。定理 6.6.11 蕴涵周期映射是满射。为理解 ALH$^*$ 引力瞬子的空间，人们会想刻画周期区域的像。这将在即将进行的工作 [47] 中研究。

另一个问题来自如下观察：ALH$^*$ 引力瞬子经超 Kähler 旋转后的不同紧化，启示了其他类型引力瞬子的不同紧化之可能。已经证明 ALG 与 ALG$^*$ 引力瞬子在适当的超 Kähler 旋转后可以紧化为有理椭圆曲面 [6][14]。

\begin{conjecture}{6.6.15}
ALG 或 ALG$^*$ 引力瞬子（在适当的超 Kähler 旋转下）是否容许紧化为有理椭圆曲面之外的有理曲面（通过添加一个典范圈得到）？例如，该猜想对 ALG 引力瞬子在适当的超 Kähler 旋转后成立 [13]。
\end{conjecture}

\subsection{辛同胚}

设 $Y$ 是次数为 $d$ 的 del Pezzo 曲面，$D\subseteq Y$ 是光滑反典范除子，$(\check Y,\check D)$ 是定理 6.2.3 中的有理椭圆曲面配对。则由 (6.2.2)，配对 $(\check Y,\check D)$ 的任一保持 $\check Y$ 上沿 $\check D$ 有单极点的亚纯体积形式的自同构，都给出 $X=Y\setminus D$ 关于 Tian--Yau 度量的一个辛同胚 [11, 命题 5.10]。此外，作为定理 6.5.1 的直接推论，我们有如下观察：

\begin{proposition}{6.6.16}
在上述设定下，若 $\phi:(\check Y,\check D)\cong(\check Y,\check D)$ 是配对的同构，使得
% [原文 p.26]
\begin{enumerate}
\item $\phi$ 保持 $\check Y$ 上沿 $\check D$ 有单极点的亚纯体积形式，且
\item $\phi$ 保持 $\check X=\check Y\setminus\check D$ 从超 Kähler 旋转所得的 Kähler 类，
\end{enumerate}
则超 Kähler 旋转之前的同一个底层映射（仍记作 $\phi$）是配对 $(Y,D)$ 的自同构。
\end{proposition}

特别地，在 $Y=\PP^2$ 的情形我们恢复一个经典定理。

\begin{theorem}{6.6.17}
给定 $\PP^2$ 中任一光滑三次曲线 $D$，存在 $\operatorname{Aut}(\PP^2)$ 的子群 $\Z_3\oplus\Z_3$ 保持 $D$。
\end{theorem}

\begin{remark}{6.6.18}
命题 6.6.16 到 Looijenga 对的推广曾通过镜像对称在 [34, 42] 中得证。
\end{remark}
% [Tutor Note: 本章正文止于本条注 6.6.18，其后即 p.26 的 REFERENCES；全篇（p.1--p.26）无“致谢”一节，故本块不写 \section*{致谢}。]
% [Tutor Note: 第 6.3 节标题原文印作 “6.3. Sympectomorphisms.”（Symplectomorphisms 的拼写笔误，页 25 可辨），译文按“辛同胚”处理；本节以 $\phi$ 表示配对自同构，与前几节用 $\Phi$、$\varphi$ 的对象不同，抄写时勿混。]
